\documentclass[a4paper,11pt]{article}
\usepackage[utf8]{inputenc}
\usepackage[T1]{fontenc}
\usepackage[french]{babel}
\usepackage[left=2.5cm,right=2.5cm,top=1.8cm,bottom=1.8cm]{geometry}
\usepackage{hyperref}
\usepackage{xcolor}
\usepackage{fontawesome5}
\usepackage{lmodern}

\definecolor{primary}{RGB}{0, 82, 147}
\hypersetup{colorlinks=true, linkcolor=primary, urlcolor=primary}

\begin{document}

\noindent
\textbf{Loïc Aron MBASSI EWOLO} \\
Élève-Ingénieur ENSEA / Polytechnique Yaoundé \\
Cergy, France \\
\faPhone\ (+33) 7 59 52 33 98 \quad \faEnvelope\ \href{mailto:loicaron.mbassiewolo@ensea.fr}{loicaron.mbassiewolo@ensea.fr}

\vspace{0.6cm}

\noindent
\textbf{CEA DAM Île-de-France} \\
À l'attention de Dr. Martin COLLET \\
Bruyères-le-Châtel - 91297 Arpajon

\hfill Cergy, le 23 janvier 2026

\vspace{0.5cm}

\noindent\textbf{Objet : Stage - Développement d'un système de compteur d'impulsions (ELSA)}

\vspace{0.4cm}

Monsieur,

\vspace{0.3cm}

Je suis en deuxième année à l'ENSEA, en double diplôme avec Polytechnique Yaoundé. Je cherche un stage de 4 à 6 mois, et votre offre sur le système de compteur d'impulsions m'intéresse pour une raison simple : c'est exactement le type de projet où je peux être utile.

\vspace{0.25cm}

À l'ENSEA, je suis la filière électronique et systèmes embarqués. J'ai des cours d'électronique analogique et numérique, de microcontrôleurs, et d'instrumentation. Le contexte technique de l'offre (signaux créneau, synchronisation, conception de prototype) m'est donc familier, même si je n'ai jamais travaillé sur un accélérateur d'électrons.

\vspace{0.25cm}

Ce qui me parle dans cette offre, c'est la partie conception de prototype à base de microcontrôleur. J'ai déjà fait ça. Au laboratoire IoT de mon école, j'ai bossé sur des systèmes embarqués avec Arduino et Raspberry Pi. Plus récemment, pour un projet de traduction de langue des signes, j'ai participé à la conception d'un gant instrumenté : soudure de capteurs IMU et flex sensors, acquisition de signaux en temps réel, traitement embarqué sur STM32. Ce n'était pas un compteur d'impulsions, mais le processus (analyser un besoin, prototyper, itérer) est le même.

\vspace{0.25cm}

Sur la programmation, je code en Python et en C. J'ai aussi utilisé MATLAB/Simulink pour un projet de contrôle de drone à l'ENSEA. L'algorithme de gestion des impulsions que vous décrivez me semble abordable avec ces outils. Et si les méthodes d'IA peuvent être explorées, j'ai une expérience de recherche au Mila (Québec) sur les réseaux de neurones, même si je ne suis pas sûr que ce soit pertinent pour ce stage.

\vspace{0.25cm}

Concernant la rédaction du cahier des charges, j'ai déjà documenté une API financière chez CSC (mission freelance) pour qu'elle soit reprise par d'autres développeurs. Cette expérience m'a appris à structurer un document technique de manière claire.

\vspace{0.25cm}

Je suis basé à Cergy, donc Bruyères-le-Châtel est accessible. Je suis disponible à partir d'avril 2026.

\vspace{0.4cm}

Cordialement,

\vspace{0.8cm}

\textbf{Loïc Aron MBASSI EWOLO}

\end{document}
